\documentclass[11pt]{article}

\usepackage[margin=1in]{geometry}
\usepackage{amsmath}
\usepackage{amssymb}
\usepackage{booktabs}
\usepackage{graphicx}
\usepackage{hyperref}

\title{Supporting Information\\[0.3em]
  \large Definitional Instability in Kinase Inhibitor Selectivity:
  An Empirical Characterization and Motivated Properties
  for a Well-Formed Measure}
\author{Polina Vinogradova}
\date{}

\begin{document}
\maketitle

This document reports, in full, four robustness checks summarized in the main
text. All numbers are reproducible from the code repository
(\url{https://github.com/polinavino/kinase-selectivity-definitions}); the
script generating each table is named in its caption.

\section*{Table S1: Panel-size recoverability under alternative correlation
bars}

The main text calls a panel size ``recovering'' the full-panel ranking once
mean Spearman $\rho$ against the full panel exceeds $0.90$, a convention
borrowed from reliability research's ``excellent agreement'' bar rather than
derived from the data (Methods, Results). Table~S1 reports $p^*$ (the
smallest panel size reaching the bar) at four alternative bars, for both
sparse-coverage datasets.

\begin{table}[h]
\centering
\caption{Panel-size $p^*$ (kinases) at four correlation bars.
(\texttt{pstar\_bar\_sensitivity.py})}
\begin{tabular}{lcccc}
\toprule
Klaeger (343 kinases) & entropy & gini & s\_score & ratio \\
\midrule
bar = 0.80 & 80  & 260 & 50  & 320 \\
bar = 0.85 & 80  & 260 & 80  & 320 \\
bar = 0.90 & 110 & 290 & 110 & 343 \\
bar = 0.95 & 140 & 320 & 200 & 343 \\
\bottomrule
\end{tabular}

\vspace{1em}
\begin{tabular}{lcccc}
\toprule
Metz (172 kinases) & entropy & gini & s\_score & ratio \\
\midrule
bar = 0.80 & 50  & 50  & 50 & 170 \\
bar = 0.85 & 50  & 50  & 50 & 170 \\
bar = 0.90 & 80  & 80  & 50 & 170 \\
bar = 0.95 & 110 & 110 & 80 & 170 \\
\bottomrule
\end{tabular}
\end{table}

The ratio remains the slowest to recover and the Gini coefficient the
second-slowest for every bar tested; only the relative order of entropy and
the S-score is bar-dependent, and only within $[0.80, 0.95]$.

\clearpage
\section*{Table S2: Bosc et al.\ window/ranking score and the CATDS proxy}

Bosc et al.'s window score (WS) and ranking score (RS) and a CATDS-style
proxy (Discussion) were computed on all four datasets to check whether they
behave as additional members of the ratio family (WS, RS) or the
distribution-based family (CATDS), as claimed in the Discussion.
(\texttt{bosc\_catds\_comparison.py})

\begin{table}[h]
\centering
\caption{Spearman correlation of WS/RS against the ratio family, and of the
CATDS proxy against the ratio, entropy, and Gini families.}
\begin{tabular}{lccccc}
\toprule
Dataset & $n_\mathrm{cpd}$ & WS vs.\ ratio & RS vs.\ ratio & CATDS vs.\ ratio &
CATDS vs.\ entropy / Gini \\
\midrule
Davis         & 68  & 0.804 & 0.878 & 0.509 & 0.955 / 0.960 \\
Klaeger       & 222 & 0.909 & 0.892 & 0.681 & 0.850 / 0.985 \\
Anastassiadis & 178 & 0.959 & 0.891 & 0.696 & 0.948 / 0.965 \\
Metz          & 704 & 0.910 & 0.899 & 0.447 & 0.883 / 0.906 \\
\bottomrule
\end{tabular}
\end{table}

WS and RS correlate strongly with the ratio family on every dataset ($r =
0.80$--$0.96$), supporting their classification alongside $R(k)$. The CATDS
proxy instead correlates more strongly with entropy and Gini than with the
ratio on every dataset, supporting the Discussion's claim that CATDS is a
distributional construction despite its historical classification as a ratio
measure.

\clearpage
\section*{Table S3: Cross-definition correlations under the original metric
operationalizations}

Methods reports summary statistics for this check; the full $4\times4$
correlation matrices are given here.
(\texttt{original\_operationalization\_robustness.py})

\begin{table}[h]
\centering
\caption{Davis ($n=68$): manuscript S-score/Gini/ratio vs.\ literal
Ka-weighted entropy (complete matrix).}
\begin{tabular}{lcccc}
\toprule
 & S-score & entropy(literal) & Gini & ratio \\
\midrule
S-score & 1.000 & 0.417 & 0.952 & 0.456 \\
entropy(literal) & 0.417 & 1.000 & 0.424 & 0.336 \\
Gini & 0.952 & 0.424 & 1.000 & 0.343 \\
ratio & 0.456 & 0.336 & 0.343 & 1.000 \\
\bottomrule
\end{tabular}

\vspace{1em}
\caption*{Klaeger ($n=222$): manuscript S-score/Gini/ratio vs.\ literal
Ka-weighted entropy (measured high-confidence entries only, per compound).}
\begin{tabular}{lcccc}
\toprule
 & S-score & entropy(literal) & Gini & ratio \\
\midrule
S-score & 1.000 & 0.570 & 0.722 & 0.223 \\
entropy(literal) & 0.570 & 1.000 & 0.688 & 0.789 \\
Gini & 0.722 & 0.688 & 1.000 & 0.616 \\
ratio & 0.223 & 0.789 & 0.616 & 1.000 \\
\bottomrule
\end{tabular}

\vspace{1em}
\caption*{Metz ($n=704$): manuscript S-score/Gini/ratio vs.\ literal
Ka-weighted entropy (measured entries only, per compound).}
\begin{tabular}{lcccc}
\toprule
 & S-score & entropy(literal) & Gini & ratio \\
\midrule
S-score & 1.000 & 0.239 & 0.879 & 0.211 \\
entropy(literal) & 0.239 & 1.000 & 0.416 & 0.758 \\
Gini & 0.879 & 0.416 & 1.000 & 0.229 \\
ratio & 0.211 & 0.758 & 0.229 & 1.000 \\
\bottomrule
\end{tabular}

\vspace{1em}
\caption*{Anastassiadis ($n=178$): manuscript S-score/entropy/ratio vs.\
native \%inhibition Gini (no baseline subtraction).}
\begin{tabular}{lcccc}
\toprule
 & S-score & entropy & Gini(native) & ratio \\
\midrule
S-score & 1.000 & 0.880 & 0.345 & 0.489 \\
entropy & 0.880 & 1.000 & 0.632 & 0.529 \\
Gini(native) & 0.345 & 0.632 & 1.000 & 0.314 \\
ratio & 0.489 & 0.529 & 0.314 & 1.000 \\
\bottomrule
\end{tabular}
\end{table}

On Klaeger and Metz, literal entropy correlates more strongly with the ratio
than with either S-score or Gini -- under this operationalization, entropy
would not be classified in the distribution-based family at all on either
sparse-coverage dataset. On Davis, literal entropy decouples from both
families (consistent with Figure~1 of the main text). The native-Gini check on
Anastassiadis is the one case that preserves the qualitative two-family
pattern, albeit more weakly than the manuscript's baseline-swept form.

\clearpage
\section*{Table S4: Candidate measure cross-dataset validation}

Orientation (correlation with the number of active kinases, more negative is
more correctly oriented toward ``fewer active targets = more selective''), D3
robustness under the emphasis parameter $\beta$ with the gate $x_0$ fixed
(worst-case pairwise Spearman $\rho$), D3 robustness when $x_0$ itself is
swept as a free parameter (``gate swept'', the equivalent comparison to the
entropy's baseline sensitivity; see Figure~6 of the main text), and the
percentage of compounds for which D4 (monotonicity under a weak additional
interaction) holds exactly.
(\texttt{candidate\_validation.py})

\begin{table}[h]
\centering
\caption{Candidate measure: orientation, D3 (two variants), and D4, across
all four datasets.}
\begin{tabular}{lccccc}
\toprule
Dataset & $n_\mathrm{cpd}$ & orientation $r$ & D3, gate fixed & D3, gate swept
& D4 (\% compounds) \\
\midrule
Davis         & 68  & $-0.969$ & $0.998$ & $-0.838$ & 100.0 \\
Klaeger       & 222 & $-0.936$ & $0.999$ & $-0.932$ & 100.0 \\
Anastassiadis & 178 & $-0.936$ & $0.997$ & $+0.913$ & 100.0 \\
Metz          & 704 & $-0.791$ & $0.991$ & $-0.330$ & 100.0 \\
\bottomrule
\end{tabular}
\end{table}

The ``gate swept'' column confirms, on three of four datasets, that the
candidate's D3 satisfaction rests on fixing the gate $x_0$ by declaration
rather than on the measure being inherently insensitive to its placement
(worst-case $\rho$ is far from $1$, and on Klaeger and Davis is strongly
negative). Anastassiadis is an exception ($\rho = +0.913$): on that dataset
the candidate's ranking happens to be comparatively insensitive to the gate's
placement as well, for reasons we do not have a mechanistic account of.

Smoothing-width $T$ was also swept on Klaeger ($T \in \{0.5, 0.75, 1.0, 1.5,
2.0\}$): orientation and D3-emphasis-robustness are essentially unchanged
across this range ($r \approx -0.94$, $\rho \geq 0.989$), and the worst-case
pairwise ranking agreement across all five values of $T$ is $\rho = 0.999$.
Removing the gate inverts the measure's orientation entirely (gated $r =
-0.936$ vs.\ ungated $r = +0.943$ against $n_\mathrm{active}$ on Klaeger),
which is why the gate, not merely the smooth emphasis, is required for D1 and
correct orientation jointly.

\end{document}
